% Fuente: Auxiliar 3, sección «Método gráfico y soluciones básicas factibles».
Un granjero posee 45 hectáreas de tierra. Va a plantar cada hectárea con trigo o maíz. Cada hectárea sembrada con trigo genera una ganancia de 200 [\$/hectárea]; cada hectárea con maíz genera una ganancia de 300 [\$/hectárea]. La mano de obra y el fertilizante utilizados por hectárea están dados en la Tabla \ref{ex:3:3:requerimientos}. Se dispone de 100 trabajadores y 120 toneladas de fertilizante.

\begin{table}[h!]
\centering
\begin{tabular}{|c|c|c|}
\hline
\textbf{} & \textbf{Trigo} & \textbf{Maíz} \\
\hline
Mano de obra & 3 trabajadores & 2 trabajadores \\
\hline
Fertilizante & 2 toneladas & 4 toneladas \\
\hline
\end{tabular}
\caption{Requerimientos por hectárea}
\label{ex:3:3:requerimientos}
\end{table}

\begin{enumerate}
    \item[a)] Modele el problema como un LP y luego obtenga la forma estándar.
    \item[b)] Resuelva gráficamente el problema. Encuentre las soluciones básicas factibles.
    \item[c)] Recientemente, el granjero firmó un contrato para suministrar maíz a una empresa procesadora, lo que requiere que se planten al menos 25 hectáreas de maíz. ¿Cómo cambia el problema con respecto a la solución anterior?
\end{enumerate}

